\section{Dynamical semigroups and the exponential form}
\label{sec:semigroup_defs}
%% ---------------------------------------------------------

\begin{definition}[Dynamical semigroup]\label{def:dyn_semigroup_ch12}
    \lean{IsDynSemigroup}\leanok
    A family of linear maps
    $T : \R \to (\MN{D} \to_{\ell} \MN{D})$
    is a \emph{dynamical semigroup} if
    \begin{enumerate}
        \item (semigroup law) $T_{t+s} = T_t \circ T_s$
              for all $t,s \ge 0$; and
        \item (initial condition) $T_0 = \id$.
    \end{enumerate}
    This is \cite[Equation~(7.1)]{Wolf2012Quantum}.
\end{definition}

\begin{definition}[Continuous dynamical semigroup]\label{def:continuous_dyn_semigroup_ch12}
    \lean{IsContinuousDynSemigroup}\leanok
    \uses{def:dyn_semigroup_ch12}
    A dynamical semigroup $T$ is \emph{norm-continuous} if the map
    $t \mapsto T_t$ is continuous in the operator norm topology.
    In finite dimension this coincides with strong continuity.
\end{definition}

\begin{definition}[Exponential semigroup]\label{def:exp_semigroup_ch12}
    \lean{expSemigroup, expSemigroupCLM}\leanok
    Given a generator $L \in \End_{\C}(\MN{D})$, the
    \emph{exponential semigroup} is
    \begin{align}
        T_t
        &:= e^{tL}
          = \sum_{k=0}^{\infty} \frac{(tL)^k}{k!}.
        \label{eq:semigroup_exp_def}
    \end{align}
\end{definition}

\begin{theorem}[Semigroup law for the exponential]\label{thm:expSemigroupCLM_add_ch12}
    \lean{expSemigroupCLM_add, expSemigroup_comp}\leanok
    \uses{def:exp_semigroup_ch12}
    For all $t,s \in \R$, $e^{(t+s)L} = e^{tL} \cdot e^{sL}$.
\end{theorem}

\begin{proof}
    \leanok
    Since $(t \cdot L)$ commutes with $(s \cdot L)$, this follows from
    the exponential addition formula for commuting elements in a
    Banach algebra.
\end{proof}

\begin{theorem}[Exponential semigroup at time $0$]
    \label{thm:expSemigroupCLM_zero_ch12}
    \lean{expSemigroupCLM_zero}\leanok
    \uses{def:exp_semigroup_ch12}
    The continuous-linear-map exponential satisfies $e^{0L} = \Id$.
\end{theorem}

\begin{proof}
    \leanok
    The exponential of the zero endomorphism is the identity.
\end{proof}

\begin{theorem}[Continuity of the exponential semigroup]\label{thm:expSemigroupCLM_continuous_ch12}
    \lean{expSemigroupCLM_continuous, expSemigroup_isContinuousDynSemigroup}\leanok
    \uses{def:exp_semigroup_ch12}
    The map $t \mapsto e^{tL}$ is continuous in the operator norm topology.
\end{theorem}

\begin{proof}
    \leanok
    The exponential is analytic (convergence radius $= \infty$), hence
    continuous.
\end{proof}

\begin{theorem}[Derivative of the exponential semigroup]\label{thm:hasDerivAt_expSemigroupCLM_ch12}
    \lean{hasDerivAt_expSemigroupCLM}\leanok
    \uses{def:exp_semigroup_ch12}
    For all $t \in \R$, one has
    $\frac{d}{dt}e^{tL} = e^{tL} \cdot L$.
    This is \cite[Equation~(7.2)]{Wolf2012Quantum}.
\end{theorem}

\begin{proof}
    \leanok
    Apply the chain rule to the composition
    $\R \xrightarrow{t \mapsto t \cdot L} \End(\MN{D})
        \xrightarrow{\exp} \End(\MN{D})$.
\end{proof}

\begin{theorem}[Derivative at the origin]
    \label{thm:hasDerivAt_expSemigroupCLM_zero_ch12}
    \lean{hasDerivAt_expSemigroupCLM_zero}\leanok
    \uses{thm:hasDerivAt_expSemigroupCLM_ch12, thm:expSemigroupCLM_zero_ch12}
    One has $\left.\frac{d}{dt}e^{tL}\right|_{t=0} = L$.
\end{theorem}

\begin{proof}
    \leanok
    This is Theorem~\ref{thm:hasDerivAt_expSemigroupCLM_ch12} at $t=0$
    together with Theorem~\ref{thm:expSemigroupCLM_zero_ch12}.
\end{proof}

\begin{theorem}[Generator uniqueness]\label{thm:generator_unique_ch12}
    \lean{generator_unique}\leanok
    \uses{def:exp_semigroup_ch12}
    If $e^{tL} = e^{tL'}$ for all $t \ge 0$, then $L = L'$.
\end{theorem}

\begin{proof}
    \leanok
    Both semigroups agree on $[0,\infty)$, so their derivatives within
    $[0,\infty)$ at $t=0$ agree.
    Since $[0,\infty)$ has unique differentials at~$0$, $L = L'$.
\end{proof}

\begin{theorem}[CLM and linear-map forms coincide]
    \label{thm:expSemigroup_toCLM_ch12}
    \lean{expSemigroup_toCLM}\leanok
    \uses{def:exp_semigroup_ch12}
    The linear-map exponential semigroup and the continuous-linear-map
    exponential semigroup agree under the canonical identification of
    endomorphisms with continuous linear endomorphisms.
\end{theorem}

\begin{proof}
    \leanok
    This is a direct unpacking of the definition of the linear-map
    exponential semigroup.
\end{proof}

\begin{theorem}[Pointwise derivative formula]
    \label{thm:hasDerivAt_expSemigroup_apply_ch12}
    \lean{hasDerivAt_expSemigroup_apply}\leanok
    \uses{thm:hasDerivAt_expSemigroupCLM_ch12, thm:expSemigroup_toCLM_ch12}
    For every $X \in \MN{D}$ and every $t \in \R$,
    $\frac{d}{dt}(e^{tL}(X)) = e^{tL}(L(X))$.
\end{theorem}

\begin{proof}
    \leanok
    Evaluate the derivative of the continuous-linear-map semigroup at $X$
    and transport the statement through
    Theorem~\ref{thm:expSemigroup_toCLM_ch12}.
\end{proof}

\begin{theorem}[Linear-map exponential semigroup at time $0$]
    \label{thm:expSemigroup_zero_ch12}
    \lean{expSemigroup_zero}\leanok
    \uses{thm:expSemigroup_toCLM_ch12, thm:expSemigroupCLM_zero_ch12}
    The linear-map exponential semigroup satisfies $e^{0L} = \Id$.
\end{theorem}

\begin{proof}
    \leanok
    Transport Theorem~\ref{thm:expSemigroupCLM_zero_ch12} through
    Theorem~\ref{thm:expSemigroup_toCLM_ch12}.
\end{proof}

\begin{theorem}[Exponential semigroup as a dynamical semigroup]
    \label{thm:expSemigroup_isDynSemigroup_ch12}
    \lean{expSemigroup_isDynSemigroup}\leanok
    \uses{def:dyn_semigroup_ch12, thm:expSemigroup_toCLM_ch12,
        thm:expSemigroup_zero_ch12, thm:expSemigroupCLM_add_ch12}
    The family $t \mapsto e^{tL}$ satisfies the semigroup law and the initial
    condition, hence is a dynamical semigroup.
\end{theorem}

\begin{proof}
    \leanok
    Transport the continuous-linear-map semigroup law from
    Theorem~\ref{thm:expSemigroupCLM_add_ch12} through
    Theorem~\ref{thm:expSemigroup_toCLM_ch12}; the initial condition is
    Theorem~\ref{thm:expSemigroup_zero_ch12}.
\end{proof}

\begin{theorem}[Every continuous semigroup is exponential]
    \label{thm:continuousDynSemigroup_eq_exp_ch12}
    \lean{continuousDynSemigroup_eq_exp}\leanok
    \uses{def:continuous_dyn_semigroup_ch12, def:exp_semigroup_ch12}
    Every norm-continuous dynamical semigroup $T$ on $\MN{D}$ is
    of the form $T_t = e^{tL}$ for some generator
    $L \in \End_{\C}(\MN{D})$ and every $t \ge 0$.
    This is \cite[Proposition~7.1]{Wolf2012Quantum}.
\end{theorem}

\begin{proof}
    \leanok
    Since $T_0 = \Id$ and $T$ is continuous, the integral
    $M_\varepsilon = \int_0^\varepsilon T_s\,ds$ is invertible for small
    $\varepsilon > 0$.
    Using the semigroup property,
    $T_t = M_\varepsilon^{-1}(M_{t+\varepsilon} - M_t)$
    is differentiable for $t \ge 0$.
    ODE uniqueness then gives $T_t = e^{tL}$ for all $t \ge 0$.
\end{proof}

%% ---------------------------------------------------------
\section{Perturbation theory}
\label{sec:perturbation}
%% ---------------------------------------------------------

\begin{theorem}[Duhamel formula]\label{thm:duhamel_formula_ch12}
    \lean{duhamel_formula}\leanok
    \uses{def:exp_semigroup_ch12, thm:hasDerivAt_expSemigroupCLM_ch12}
    Let $\Delta = L' - L$. For $t \ge 0$,
    \begin{align}
        e^{tL'} - e^{tL}
        &= \int_0^t e^{(t-s)L} \Delta\,e^{sL'}\,ds.
        \label{eq:semigroup_duhamel}
    \end{align}
    This is \cite[Lemma~7.1]{Wolf2012Quantum}.
\end{theorem}

\begin{proof}
    \leanok
    Define $f(s) = e^{(t-s)L}e^{sL'}$.
    Then $f'(s) = e^{(t-s)L}\Delta\,e^{sL'}$ by
    Theorem~\ref{thm:hasDerivAt_expSemigroupCLM_ch12}, and
    $e^{tL'} - e^{tL} = f(t) - f(0) = \int_0^t f'(s)\,ds$,
    which gives (\ref{eq:semigroup_duhamel}).
\end{proof}

\begin{theorem}[Perturbation bound]\label{thm:perturbation_bound_ch12}
    \lean{perturbation_bound}\leanok
    \uses{thm:duhamel_formula_ch12}
    For $t \ge 0$,
    \begin{align}
        \bigl\|e^{tL'} - e^{tL}\bigr\|
        &\le t\,\|\Delta\|
        \cdot \sup_{s \in [0,t]}\|e^{sL}\|
        \cdot \sup_{s \in [0,t]}\|e^{sL'}\|,
        \label{eq:semigroup_perturbation_bound}
    \end{align}
    where $\Delta = L' - L$.
    This is \cite[Corollary~7.1]{Wolf2012Quantum}.
\end{theorem}

\begin{proof}
    \leanok
    Apply the triangle inequality and submultiplicativity of the
    operator norm to the Duhamel integral
    (\ref{eq:semigroup_duhamel}).
\end{proof}

\begin{corollary}[Perturbation bound under unit-norm control]
    \label{cor:perturbation_bound_unit_norm_ch12}
    \lean{perturbation_bound_unit_norm}\leanok
    \uses{thm:perturbation_bound_ch12}
    For $t \ge 0$, if $\|e^{sL}\| \le 1$ and $\|e^{sL'}\| \le 1$ for every
    $s \in [0,t]$, then
    $\|e^{tL'} - e^{tL}\| \le t\,\|L' - L\|$.
\end{corollary}

\begin{proof}
    \leanok
    Apply (\ref{eq:semigroup_perturbation_bound}) and bound both suprema
    by $1$.
\end{proof}

\begin{lemma}[Dyson-series summability from factorial majorant]
    \label{lem:summable_dyson_factorial_ch12}
    \lean{summable_dysonTerm_of_factorial_bound}\leanok
    Assume a factorial majorant of Dyson--Phillips iterates:
    \begin{align}
        \|D_n(t)\|
        &\le M\,\frac{(t\,\|L'-L\|\,M)^n}{n!}.
        \label{eq:semigroup_dyson_majorant}
    \end{align}
    Then the Dyson series is summable in operator norm.
\end{lemma}

\begin{proof}
    \leanok
    This is the Weierstrass M-test step, using the scalar convergence of
    $\sum_n a^n/n!$ and lifting via norm-bounded summability.
\end{proof}

\begin{theorem}[Factorial norm bound for Dyson iterates]
    \label{thm:norm_dysonTerm_le_ch12}
    \lean{norm_dysonTerm_le}\leanok
    \uses{def:exp_semigroup_ch12}
    For $s \in [0,t]$ and $M = \sup_{u \in [0,t]} \|e^{uL}\|$,
    \begin{align}
        \|\widetilde{T}^{(n)}_s\|
        &\le M\,\frac{(s\,\|\Delta\|\,M)^n}{n!}.
        \label{eq:semigroup_dyson_iterate}
    \end{align}
\end{theorem}

\begin{proof}
    \leanok
    \uses{def:exp_semigroup_ch12}
    Induction on $n$. The base case uses the semigroup supremum bound.
    The inductive step bounds the integrand pointwise by
    $M\|\Delta\| \cdot M(u\|\Delta\|M)^n/n!$ and integrates
    $\int_0^s u^n\,du = s^{n+1}/(n+1)$, recovering the factorial.
\end{proof}

\begin{corollary}[Dyson series summability]
    \label{cor:summable_dysonTerm_ch12}
    \lean{summable_dysonTerm}\leanok
    \uses{thm:norm_dysonTerm_le_ch12}
    The Dyson--Phillips series $\sum_n \widetilde{T}^{(n)}_t$
    converges in operator norm for every $t \ge 0$.
\end{corollary}

\begin{proof}
    \leanok
    \uses{thm:norm_dysonTerm_le_ch12, lem:summable_dyson_factorial_ch12}
    The bound (\ref{eq:semigroup_dyson_iterate}) gives a factorial majorant.
    Lemma~\ref{lem:summable_dyson_factorial_ch12} then applies the Weierstrass
    M-test.
\end{proof}

\begin{corollary}[Factorial norm bound at the final time]
    \label{cor:norm_dysonTerm_le_at_ch12}
    \lean{norm_dysonTerm_le_at}\leanok
    \uses{thm:norm_dysonTerm_le_ch12}
    For $t \ge 0$, setting $s=t$ in
    (\ref{eq:semigroup_dyson_iterate}) gives
    \begin{align}
        \|\widetilde{T}^{(n)}_t\|
        &\le M\,\frac{(t\,\|\Delta\|\,M)^n}{n!},
        \label{eq:semigroup_dyson_final}
    \end{align}
    where $M = \sup_{u \in [0,t]} \|e^{uL}\|$.
\end{corollary}

\begin{proof}
    \leanok
    This is the special case $s=t$ of
    (\ref{eq:semigroup_dyson_iterate}).
\end{proof}

\begin{theorem}[Continuity of Dyson iterates]
    \label{thm:dysonTerm_continuous_ch12}
    \lean{dysonTerm_continuous}\leanok
    \uses{def:exp_semigroup_ch12}
    Each Dyson iterate $t \mapsto \widetilde{T}^{(n)}_t$ is continuous
    in the time parameter.
\end{theorem}

\begin{proof}
    \leanok
    \uses{def:exp_semigroup_ch12}
    By induction on $n$. The base case is continuity of the matrix
    exponential. For the inductive step, the parametric integral
    $t\mapsto \int_0^t e^{(t-s)L}\Delta\,\widetilde{T}^{(n)}_s\,ds$
    is continuous by the continuous-parameter primitive theorem;
    the pasting lemma extends this to all of $\R$ since
    $\widetilde{T}^{(n+1)}_t = 0$ for $t \le 0$.
\end{proof}

\begin{theorem}[Factorial bound on the Dyson remainder]
    \label{thm:norm_dysonRemainder_le_ch12}
    \lean{norm_dysonRemainder_le}\leanok
    \uses{thm:duhamel_formula_ch12, thm:norm_dysonTerm_le_ch12}
    For $s \in [0,t]$, $M = \sup_{u \in [0,t]}\|e^{uL}\|$,
    and $M' = \sup_{u \in [0,t]}\|e^{uL'}\|$,
    \begin{align}
        \bigl\|T'_s - \textstyle\sum_{n < N}\widetilde{T}^{(n)}_s\bigr\|
        &\le M'\,\frac{(s\,\|\Delta\|\,M)^N}{N!}.
        \label{eq:semigroup_dyson_remainder}
    \end{align}
\end{theorem}

\begin{proof}
    \leanok
    \uses{thm:duhamel_formula_ch12, thm:norm_dysonTerm_le_ch12,
        thm:dysonTerm_continuous_ch12}
    Induction on $N$. The base case is the semigroup supremum bound.
    The inductive step uses the telescoping remainder identity
    $R_{N+1}(s) = \int_0^s e^{(s-u)L}\Delta\,R_N(u)\,du$
    derived from the Duhamel formula (\ref{eq:semigroup_duhamel}), bounds
    the integrand pointwise, and integrates to recover the factorial.
\end{proof}

\begin{theorem}[Dyson series identity {\cite[Equation~(7.13)]{Wolf2012Quantum}}]
    \label{thm:dyson_series_eq_ch12}
    \lean{dyson_series_eq}\leanok
    \uses{thm:norm_dysonRemainder_le_ch12, cor:summable_dysonTerm_ch12}
    For $t \ge 0$, the Dyson--Phillips series equals the perturbed
    semigroup:
    \begin{align}
        \sum_{n=0}^{\infty} \widetilde{T}^{(n)}_t
        &= e^{tL'}.
        \label{eq:semigroup_dyson_series}
    \end{align}
\end{theorem}

\begin{proof}
    \leanok
    \uses{thm:norm_dysonRemainder_le_ch12, cor:summable_dysonTerm_ch12}
    By (\ref{eq:semigroup_dyson_remainder}), the partial-sum remainder
    $\|e^{tL'}-\sum_{n<N}\widetilde{T}^{(n)}_t\|
        \le M'(t\|\Delta\|M)^N/N!$
    tends to zero as $N \to \infty$, since $c^N/N! \to 0$ for any
    constant $c$. Together with the summability of the series
    (Corollary~\ref{cor:summable_dysonTerm_ch12}), this identifies
    the sum via uniqueness of limits.
\end{proof}

\begin{corollary}[Tsum form of the Dyson series]
    \label{cor:tsum_dysonTerm_eq_ch12}
    \lean{tsum_dysonTerm_eq}\leanok
    \uses{thm:dyson_series_eq_ch12}
    One has $\sum_{n=0}^{\infty}\widetilde{T}^{(n)}_t = e^{tL'}$
    in tsum form.
\end{corollary}

\begin{proof}
    \leanok
    \uses{thm:dyson_series_eq_ch12}
    The claim follows from (\ref{eq:semigroup_dyson_series}).
\end{proof}

%% ---------------------------------------------------------
\section{GKSL/Lindblad generators}
\label{sec:gksl_generators}
%% ---------------------------------------------------------

This section characterizes generators of semigroups of completely
positive and trace-preserving maps.
The central result is the GKSL theorem
(Theorem~\ref{thm:gksl_iff_lindblad}),
which shows that such generators have the standard Lindblad form.

\subsection{Generator decomposition and conditional complete positivity}

\begin{definition}[Generator decomposition]\label{def:generator_decomp}
    \lean{GeneratorDecomp}\leanok
    A \emph{generator decomposition} consists of a completely positive map
    $\phi : \MN{d} \to \MN{d}$ and a matrix $\kappa \in \MN{d}$, defining
    the linear map
    \begin{align}
        L(\rho)
        &= \phi(\rho) - \kappa\rho - \rho\kappa^\dagger.
        \label{eq:semigroup_generator_decomp}
    \end{align}
    This is \cite[Equation~(7.14)]{Wolf2012Quantum}.
\end{definition}

\begin{definition}[Conditional complete positivity]\label{def:is_ccp}
    \lean{IsCCP}\leanok
    \uses{def:generator_decomp}
    A linear map $L : \MN{d} \to \MN{d}$ is \emph{conditionally
        completely positive} (CCP) if it admits a generator decomposition,
    i.e.\ there exist a CP map $\phi$ and a matrix $\kappa$ such that
    $L(\rho) = \phi(\rho) - \kappa\rho - \rho\kappa^\dagger$.
    This is condition~1 of \cite[Proposition~7.2]{Wolf2012Quantum}.
\end{definition}

\begin{definition}[Trace-annihilating map]\label{def:is_trace_annihilating}
    \lean{IsTraceAnnihilating}\leanok
    A linear map $L : \MN{d} \to \MN{d}$ is \emph{trace-annihilating} if
    $\tr(L(\rho)) = 0$ for all $\rho \in \MN{d}$.
    This is the infinitesimal version of trace preservation.
\end{definition}

\begin{definition}[Trace constraint]\label{def:is_trace_constraint}
    \lean{GeneratorDecomp.isTraceConstraint}\leanok
    \uses{def:generator_decomp}
    A generator decomposition $(\phi, \kappa)$ satisfies the
    \emph{trace constraint} if $\phi^*(\Id) = \kappa + \kappa^\dagger$,
    where $\phi^*$ is the Hilbert--Schmidt adjoint.
    This is the condition in \cite[Equation~(7.20)]{Wolf2012Quantum}.
\end{definition}

\begin{theorem}[Trace constraint implies trace-annihilating]
    \label{thm:traceAnnihilating_of_traceConstraint}
    \lean{GeneratorDecomp.traceAnnihilating_of_traceConstraint}\leanok
    \uses{def:generator_decomp, def:is_trace_annihilating,
        def:is_trace_constraint}
    If $(\phi, \kappa)$ satisfies $\phi^*(\Id) = \kappa + \kappa^\dagger$,
    then $L(\rho) = \phi(\rho) - \kappa\rho - \rho\kappa^\dagger$ is
    trace-annihilating.
\end{theorem}

\begin{proof}
    \leanok
    Let $\phi(\rho) = \sum_i K_i\rho K_i^\dagger$.
    By the cyclic property of trace,
    $\tr(L(\rho))
        = \tr((\sum_i K_i^\dagger K_i)\rho)
        - \tr(\kappa\rho)
        - \tr(\kappa^\dagger\rho)
        = \tr((\kappa+\kappa^\dagger)\rho)
        - \tr(\kappa\rho)-\tr(\kappa^\dagger\rho)=0$.
\end{proof}

\subsection{The Lindblad form}

\begin{definition}[Lindblad form]\label{def:lindblad_form}
    \lean{LindbladForm}\leanok
    A \emph{Lindblad form} consists of a Hermitian matrix $H = H^\dagger$
    (the Hamiltonian) and a family of matrices $\{L_j\}_{j=0}^{r-1}$
    (the Lindblad operators), defining the linear map
    \begin{align}
        L(\rho)
        &= i[\rho,H]
        + \sum_j \Bigl(        L_j\rho L_j^\dagger
            - \tfrac{1}{2}\{L_j^\dagger L_j,\rho\}_+
        \Bigr),
        \label{eq:semigroup_lindblad_form}
    \end{align}
    where $[A,B] = AB - BA$ and $\{A,B\}_+ = AB + BA$.
    This is \cite[Equation~(7.21)]{Wolf2012Quantum}.
\end{definition}

\begin{theorem}[Lindblad form is trace-annihilating]
    \label{thm:lindbladForm_isTraceAnnihilating}
    \lean{LindbladForm.isTraceAnnihilating}\leanok
    \uses{def:lindblad_form, def:is_trace_annihilating}
    Every Lindblad form defines a trace-annihilating linear map.
\end{theorem}

\begin{proof}
    \leanok
    The Hamiltonian part satisfies
    $\tr(i[\rho,H]) = 0$ by the cyclic property.
    Each dissipator term satisfies
    $\tr(L_j\rho L_j^\dagger)
        = \tr(L_j^\dagger L_j\rho)
        = \tr(\rho L_j^\dagger L_j)$,
    so the three terms sum to zero.
\end{proof}

\begin{theorem}[Lindblad form equals generator decomposition]
    \label{thm:lindbladForm_eq_generatorDecomp}
    \lean{LindbladForm.toLinearMap_eq_generatorDecomp}\leanok
    \uses{def:lindblad_form, def:generator_decomp}
    A Lindblad form with Hamiltonian $H$ and operators $\{L_j\}$ defines
    the same linear map as the generator decomposition $(\phi,\kappa)$
    with $\phi(\rho) = \sum_j L_j\rho L_j^\dagger$ and
    $\kappa = iH + \frac{1}{2}\sum_j L_j^\dagger L_j$.
    This is \cite[Equation~(7.24)]{Wolf2012Quantum}.
\end{theorem}

\begin{proof}
    \leanok
    Compute $\kappa^\dagger = -iH+\frac{1}{2}\sum_j L_j^\dagger L_j$
    using $H^\dagger = H$ and $(A^\dagger A)^\dagger = A^\dagger A$.
    Setting $S = \sum_j L_j^\dagger L_j$, expand
    $\phi(\rho)-\kappa\rho-\rho\kappa^\dagger
        = \sum_j L_j\rho L_j^\dagger-iH\rho-\tfrac{1}{2}S\rho
        +i\rho H-\tfrac{1}{2}\rho S$,
    which is
    $i[\rho,H]+\sum_j(L_j\rho L_j^\dagger
        -\tfrac{1}{2}L_j^\dagger L_j\rho
        -\tfrac{1}{2}\rho L_j^\dagger L_j)$.
\end{proof}

\begin{theorem}[Lindblad form is CCP]
    \label{thm:lindbladForm_isCCP}
    \lean{LindbladForm.isCCP}\leanok
    \uses{def:lindblad_form, def:is_ccp}
    Every Lindblad form defines a conditionally completely positive map.
\end{theorem}

\begin{proof}
    \leanok
    \uses{thm:lindbladForm_eq_generatorDecomp}
    By Theorem~\ref{thm:lindbladForm_eq_generatorDecomp}, the Lindblad
    form equals a generator decomposition with $\phi$ CP.
\end{proof}

\begin{definition}[Commutator form of Lindblad equation]\label{def:commutator_form}
    \lean{LindbladForm.commutatorForm}\leanok
    \uses{def:lindblad_form}
    The \emph{commutator form} of the Lindblad equation writes
    the generator as
    \begin{align}
        L(\rho)
        &= i[\rho,H]
        + \tfrac{1}{2}\sum_j
        \bigl([L_j,\rho L_j^\dagger]
        + [L_j\rho,L_j^\dagger]\bigr),
        \label{eq:semigroup_commutator_form}
    \end{align}
    where $[A,B] = AB - BA$.
    This is \cite[Equation~(7.22)]{Wolf2012Quantum}.
\end{definition}

\begin{theorem}[Commutator form equals standard Lindblad form]
    \label{thm:commutatorForm_eq_toLinearMap}
    \lean{LindbladForm.commutatorForm_eq_toLinearMap}\leanok
    \uses{def:commutator_form, def:lindblad_form}
    The commutator form (Definition~\ref{def:commutator_form}) defines
    the same linear map as the standard Lindblad form
    (Definition~\ref{def:lindblad_form}).
\end{theorem}

\begin{proof}\leanok
    Expanding the commutator brackets
    $[L_j,\rho L_j^\dagger]$ and $[L_j\rho,L_j^\dagger]$
    in (\ref{eq:semigroup_commutator_form}) yields the standard dissipator
    terms in (\ref{eq:semigroup_lindblad_form}).
\end{proof}

\subsection{Characterization of conditional complete positivity}

\begin{theorem}[Conditional complete positivity implies projected Choi positivity]
    \label{thm:ccp_implies_choi_projected}
    \lean{ccp_implies_choi_projected_posSemidef}\leanok
    \uses{def:is_ccp}
    If $L$ is CCP and $P = \Id - \ketbra{\Omega}{\Omega}$, then
    \begin{align}
        P((L\otimes\id)(\ketbra{\Omega}{\Omega}))P
        &\ge 0.
        \label{eq:semigroup_choi_projected}
    \end{align}
    This is \cite[Proposition~7.2]{Wolf2012Quantum}.
\end{theorem}

\begin{proof}
    \leanok
    Write $L(\rho) = \phi(\rho)-\kappa\rho-\rho\kappa^\dagger$.
    The left- and right-multiplication terms are annihilated by $P$, so the
    projected Choi matrix reduces to the projected Choi matrix of $\phi$, which
    is positive.
\end{proof}

\begin{theorem}[Projected Choi positivity implies conditional complete positivity]
    \label{thm:choi_projected_implies_ccp}
    \lean{choi_projected_posSemidef_implies_ccp}\leanok
    \uses{def:is_ccp}
    If $L$ is Hermiticity-preserving and
    $P((L\otimes\id)(\ketbra{\Omega}{\Omega}))P \ge 0$,
    where $P = \Id - \ketbra{\Omega}{\Omega}$,
    then $L$ is CCP.
    This is \cite[Proposition~7.2]{Wolf2012Quantum}.
\end{theorem}

\begin{proof}
    \leanok
    Decompose the Hermitian Choi matrix as
    $\tau_L = Q-\ket{\psi}\bra{\Omega}-\ket{\Omega}\bra{\psi}$
    with $Q \ge 0$ supported on $\ket{\Omega}^\perp$.
    The Choi--Jamiolkowski isomorphism applied to $Q$ gives
    the CP map $\phi$, and
    $(\kappa\otimes\Id)\ket{\Omega} = \ket{\psi}$ defines $\kappa$.
\end{proof}

\subsection{Completely positive semigroups and conditionally completely positive generators}

\begin{theorem}[CP semigroups have CCP generators]
    \label{thm:cp_semigroup_implies_ccp}
    \lean{cp_semigroup_implies_ccp_generator}\leanok
    \uses{def:cp_map, def:is_ccp, def:exp_semigroup_ch12}
    If $e^{tL}$ is CP for all $t \ge 0$, then $L$ is CCP.
    This is \cite[Proposition~7.3]{Wolf2012Quantum}.
\end{theorem}

\begin{proof}
    \leanok
    \uses{thm:choi_projected_implies_ccp}
    From $(e^{tL}\otimes\id)(\ketbra{\Omega}{\Omega}) \ge 0$,
    expand at infinitesimal $t$ and project onto $P$ to get
    $P(L\otimes\id)(\ketbra{\Omega}{\Omega})P \ge 0$.
    Then Theorem~\ref{thm:choi_projected_implies_ccp} gives CCP.
\end{proof}

\begin{theorem}[CCP generators yield CP semigroups]
    \label{thm:ccp_implies_cp_semigroup}
    \lean{ccp_generator_implies_cp_semigroup}\leanok
    \uses{def:cp_map, def:is_ccp, def:exp_semigroup_ch12}
    If $L$ is CCP, then $e^{tL}$ is CP for all $t \ge 0$.
    This is \cite[Proposition~7.3]{Wolf2012Quantum}.
\end{theorem}

\begin{proof}
    \leanok
    Approximate $e^{tL}$ by the completely positive Euler steps
    \begin{align}
        \rho
        &\mapsto (\Id-h\kappa)\rho(\Id-h\kappa)^\dagger+h\phi(\rho),
        \qquad h=\frac{t}{n+1}.
        \notag
    \end{align}
    Each step is CP, hence so are its powers. These powers converge in operator
    norm to $e^{tL}$, and the cone of CP maps is closed.
\end{proof}

\begin{theorem}[CP semigroups iff CCP generators]
    \label{thm:cp_semigroup_iff_ccp}
    \lean{cp_semigroup_iff_ccp_generator}\leanok
    \uses{def:cp_map, def:is_ccp, def:exp_semigroup_ch12}
    $e^{tL}$ is CP for all $t \ge 0$ if and only if $L$ is CCP.
\end{theorem}

\begin{proof}
    \leanok
    \uses{thm:cp_semigroup_implies_ccp, thm:ccp_implies_cp_semigroup}
    Combine Theorems~\ref{thm:cp_semigroup_implies_ccp}
    and~\ref{thm:ccp_implies_cp_semigroup}.
\end{proof}

\subsection{Freedom in generator representation}

\begin{theorem}[Traceless Kraus operators exist]
    \label{thm:exists_traceless_kraus_shift}
    \lean{exists_traceless_kraus_shift}\leanok
    Assume $d > 0$. Given any Kraus operators $\{K_j\}$, there exist
    $c_j \in \C$ such that $K_j' = K_j+c_j\Id$ is traceless.
    This is part of \cite[Proposition~7.4]{Wolf2012Quantum}.
\end{theorem}

\begin{proof}
    \leanok
    Set $c_j = -\tr(K_j)/d$.
\end{proof}

\begin{theorem}[Shift invariance of generators]
    \label{thm:generator_shift_invariance}
    \lean{generator_shift_invariance}\leanok
    If $K_j' = K_j+c_j\Id$ and $\kappa'$ is adjusted per
    \cite[Equation~(7.19)]{Wolf2012Quantum},
    then $(\phi',\kappa')$ and $(\phi,\kappa)$ define the same generator.
\end{theorem}

\begin{proof}
    \leanok
    Direct algebraic expansion.
\end{proof}

\subsection{Uniqueness of the traceless Lindblad form}

\begin{definition}[Traceless Kraus operators]\label{def:has_traceless_kraus}
    \lean{LindbladForm.HasTracelessKraus}\leanok
    \uses{def:lindblad_form}
    A Lindblad form has \emph{traceless Kraus operators} if
    $\tr(L_j) = 0$ for every $j$.
\end{definition}

\begin{theorem}[Uniqueness of the completely positive part]
    \label{thm:generatorDecomp_traceless_unique_phi}
    \lean{generatorDecomp_traceless_unique_phi}\leanok
    \uses{def:lindblad_form, def:has_traceless_kraus, def:generator_decomp}
    If two Lindblad forms $F,F'$ induce the same generator and
    both have traceless Kraus operators, then their CP parts agree:
    $\phi = \phi'$.
    This is part of \cite[Proposition~7.4(2)]{Wolf2012Quantum}.
\end{theorem}

\begin{proof}\leanok
    Compare projected Choi matrices and use Choi injectivity.
\end{proof}

\begin{theorem}[Uniqueness of the drift term modulo phase]
    \label{thm:generatorDecomp_traceless_unique_kappa}
    \lean{generatorDecomp_traceless_unique_kappa_modPhase}\leanok
    \uses{def:lindblad_form, def:has_traceless_kraus, def:generator_decomp,
        thm:generatorDecomp_traceless_unique_phi}
    If two Lindblad forms $F,F'$ induce the same generator and
    both have traceless Kraus operators, then their drift matrices
    differ by an imaginary scalar: $\kappa' = \kappa+i\lambda\,\Id$
    for some $\lambda \in \R$.
    This is part of \cite[Proposition~7.4(2)]{Wolf2012Quantum}.
\end{theorem}

\begin{proof}\leanok
    \uses{thm:generatorDecomp_traceless_unique_phi}
    After identifying $\phi = \phi'$
    (Theorem~\ref{thm:generatorDecomp_traceless_unique_phi}),
    the generator equality forces
    $\Delta\kappa\cdot\rho+\rho\cdot(\Delta\kappa)^\dagger=0$
    for all~$\rho$. Setting $\rho=\Id$ gives
    $\Delta\kappa+(\Delta\kappa)^\dagger=0$
    (skew-Hermiticity), which converts the equation to
    $[\Delta\kappa,\rho]=0$ for all~$\rho$.
    The scalar commutant lemma gives $\Delta\kappa=c\cdot\Id$,
    and skew-Hermiticity forces $c=i\lambda$ for some $\lambda\in\R$.
\end{proof}

\begin{theorem}[Combined uniqueness of the traceless Lindblad form]
    \label{thm:generatorDecomp_traceless_unique}
    \lean{generatorDecomp_traceless_unique}\leanok
    \uses{thm:generatorDecomp_traceless_unique_phi,
        thm:generatorDecomp_traceless_unique_kappa}
    If two Lindblad forms $F,F'$ induce the same generator and
    both have traceless Kraus operators, then $\phi=\phi'$ and
    $\kappa'=\kappa+i\lambda\,\Id$ for some $\lambda\in\R$.
\end{theorem}

\begin{proof}\leanok
    \uses{thm:generatorDecomp_traceless_unique_phi,
        thm:generatorDecomp_traceless_unique_kappa}
    Combine Theorems~\ref{thm:generatorDecomp_traceless_unique_phi}
    and~\ref{thm:generatorDecomp_traceless_unique_kappa}.
\end{proof}

\subsection{Trace-annihilation and trace preservation}

\begin{theorem}[Trace-annihilating generators yield trace-preserving semigroups]
    \label{thm:ta_implies_tp_semigroup}
    \lean{isTracePreservingMap_expSemigroup_of_isTraceAnnihilating}\leanok
    \uses{def:is_trace_annihilating, def:trace_preserving, def:exp_semigroup_ch12}
    If $L$ is trace-annihilating, then $e^{tL}$ is trace-preserving for every
    $t \in \R$.
\end{theorem}

\begin{proof}
    \leanok
    For fixed $\rho$, the function $f(t)=\tr(e^{tL}(\rho))$ has derivative
    $f'(t)=\tr(L(e^{tL}(\rho)))=0$, so $f$ is constant and
    $\tr(e^{tL}(\rho))=f(t)=f(0)=\tr(\rho)$.
\end{proof}

\begin{theorem}[Trace-preserving semigroups have trace-annihilating generators]
    \label{thm:tp_semigroup_implies_ta}
    \lean{isTraceAnnihilating_of_isTracePreservingMap_semigroup}\leanok
    \uses{def:is_trace_annihilating, def:trace_preserving, def:exp_semigroup_ch12}
    If $e^{tL}$ is trace-preserving for every $t \ge 0$, then $L$ is
    trace-annihilating.
\end{theorem}

\begin{proof}
    \leanok
    Differentiate the identity $\tr(e^{tL}(\rho))=\tr(\rho)$ at $t=0$
    from the right.
\end{proof}

\subsection{The GKSL/Lindblad theorem}

\begin{definition}[GKSL generator]\label{def:is_gksl_generator}
    \lean{IsGKSLGenerator}\leanok
    \uses{def:channel, def:exp_semigroup_ch12}
    A linear map $L : \MN{d} \to \MN{d}$ is a \emph{GKSL generator}
    if $e^{tL}$ is a quantum channel (CPTP) for all $t \ge 0$.
\end{definition}

\begin{theorem}[A trace-constrained generator decomposition gives a GKSL generator]
    \label{thm:gksl_of_generatorDecomp}
    \lean{gksl_of_generatorDecomp_with_traceConstraint}\leanok
    \uses{def:cp_map, def:is_gksl_generator, def:generator_decomp,
        def:is_trace_constraint}
    If $L(\rho)=\phi(\rho)-\kappa\rho-\rho\kappa^\dagger$ with
    $\phi$ CP and $\phi^*(\Id)=\kappa+\kappa^\dagger$,
    then $L$ is a GKSL generator.
    This is \cite[Theorem~7.1, Equation~(7.20)]{Wolf2012Quantum}.
\end{theorem}

\begin{proof}
    \leanok
    \uses{thm:ccp_implies_cp_semigroup, thm:traceAnnihilating_of_traceConstraint,
        thm:ta_implies_tp_semigroup}
    CP follows from Theorem~\ref{thm:ccp_implies_cp_semigroup}.
    The trace constraint implies trace-annihilation by
    Theorem~\ref{thm:traceAnnihilating_of_traceConstraint}, so
    Theorem~\ref{thm:ta_implies_tp_semigroup} gives trace preservation.
\end{proof}

\begin{theorem}[GKSL generators admit a trace-constrained generator decomposition]
    \label{thm:generatorDecomp_of_gksl}
    \lean{generatorDecomp_of_gksl}\leanok
    \uses{def:cp_map, def:is_gksl_generator, def:generator_decomp, def:is_trace_constraint}
    If $L$ is a GKSL generator, then there exist a CP map $\phi$ and a matrix
    $\kappa$ such that
    \begin{align}
        L(\rho)
        &= \phi(\rho)-\kappa\rho-\rho\kappa^\dagger,
        \label{eq:semigroup_gksl_decomp}\\
        \phi^*(\Id)
        &= \kappa+\kappa^\dagger.
        \label{eq:semigroup_gksl_constraint}
    \end{align}
\end{theorem}

\begin{proof}\leanok
    \uses{thm:cp_semigroup_implies_ccp, thm:tp_semigroup_implies_ta}
    Apply Theorem~\ref{thm:cp_semigroup_implies_ccp} to obtain a CCP
    decomposition. Then apply Theorem~\ref{thm:tp_semigroup_implies_ta} to the
    trace-preserving part of the semigroup and rewrite the resulting
    trace-annihilation condition as (\ref{eq:semigroup_gksl_constraint}).
\end{proof}

\begin{theorem}[GKSL iff conditional complete positivity and trace annihilation]
    \label{thm:gksl_iff_ccp_ta}
    \lean{gksl_iff_ccp_and_traceAnnihilating}\leanok
    \uses{def:is_gksl_generator, def:is_ccp, def:is_trace_annihilating}
    $L$ is a GKSL generator if and only if $L$ is CCP and
    trace-annihilating.
\end{theorem}

\begin{proof}\leanok
    \uses{thm:cp_semigroup_implies_ccp, thm:ccp_implies_cp_semigroup,
        thm:ta_implies_tp_semigroup, thm:tp_semigroup_implies_ta}
    Forward: apply Theorem~\ref{thm:cp_semigroup_implies_ccp} to the CP part and
    Theorem~\ref{thm:tp_semigroup_implies_ta} to the TP part.
    Reverse: combine Theorem~\ref{thm:ccp_implies_cp_semigroup} with
    Theorem~\ref{thm:ta_implies_tp_semigroup}.
\end{proof}

\begin{theorem}[GKSL iff Lindblad form]
    \label{thm:gksl_iff_lindblad}
    \lean{gksl_iff_lindbladForm}\leanok
    \uses{def:is_gksl_generator, def:lindblad_form}
    $L$ is a GKSL generator if and only if
    \begin{align}
        L(\rho)
        &= i[\rho,H]
        + \sum_j \Bigl(        L_j\rho L_j^\dagger
            - \tfrac{1}{2}\{L_j^\dagger L_j,\rho\}_+
        \Bigr),
        \label{eq:semigroup_gksl_lindblad}
    \end{align}
    with $H=H^\dagger$.
    This is \cite[Theorem~7.1]{Wolf2012Quantum}.
\end{theorem}

\begin{proof}\leanok
    \uses{thm:generatorDecomp_of_gksl, thm:gksl_iff_ccp_ta,
        thm:lindbladForm_eq_generatorDecomp, thm:lindbladForm_isCCP,
        thm:lindbladForm_isTraceAnnihilating,
        thm:exists_traceless_kraus_shift}
    Forward: apply Theorem~\ref{thm:generatorDecomp_of_gksl} to write
    $L(\rho)=\phi(\rho)-\kappa\rho-\rho\kappa^\dagger$ with trace
    constraint. Then choose traceless Kraus operators by
    Theorem~\ref{thm:exists_traceless_kraus_shift} and rewrite $\kappa$ as
    $iH+\frac{1}{2}\phi^*(\Id)$;
    Theorem~\ref{thm:lindbladForm_eq_generatorDecomp} gives the Lindblad
    formula (\ref{eq:semigroup_gksl_lindblad}).
    Reverse: Theorems~\ref{thm:lindbladForm_isCCP}
    and~\ref{thm:lindbladForm_isTraceAnnihilating} show that every Lindblad form
    satisfies the right-hand side of Theorem~\ref{thm:gksl_iff_ccp_ta}.
\end{proof}

\subsection{Kossakowski matrix form}

\begin{definition}[Kossakowski form]\label{def:kossakowski_form}
    \lean{KossakowskiForm}\leanok
    A \emph{Kossakowski form} consists of $H=H^\dagger$, a finite family
    of matrices $\{F_k\}_{k=0}^{n-1}$, and a positive semidefinite matrix
    $C \in \MN{n}$, defining
    \begin{align}
        L(\rho)
        &= i[\rho,H]
        + \tfrac{1}{2}\sum_{k,l}C_{lk}
        ([F_k,\rho F_l^\dagger]+[F_k\rho,F_l^\dagger]).
        \label{eq:semigroup_kossakowski_form}
    \end{align}
    In Wolf's formula one takes $\{F_k\}$ to be a basis of traceless matrices;
    here we keep only the algebraic data needed for the conversion to
    Lindblad form.
\end{definition}

\begin{theorem}[Kossakowski form iff Lindblad form]
    \label{thm:kossakowski_iff_lindblad}
    \lean{kossakowski_iff_lindblad}\leanok
    \uses{def:kossakowski_form, def:lindblad_form}
    A linear map admits a Kossakowski form iff it admits a Lindblad form.
\end{theorem}

\begin{proof}
    \leanok
    Forward: diagonalize $C=B^\dagger B$ and set
    $L_j=\sum_k B_{jk}F_k$.
    Reverse: keep the same Hamiltonian and family $F_k=L_k$, and take
    $C=\Id$.
\end{proof}

%% ---------------------------------------------------------
